% GENERATED by scripts/manual_tables.py from the adapter registry. Do not edit:
% rerun the script. Hand-written text: blocks/interrogation/trains/
\ubblock{interrogation.trains}%
  {Trains}% what the app calls it
  {Trains of events (logISI, MaxInterval, Adamatzky's words)}% registry display name
  {control}% insert-modal tab
  {SpanSet}% input
  {SpanSet}% output
  {Adapters/interrogation\_trains.py}% code
  {interrogation/trains}% hand-written folder under blocks/
\ubblockdescription{interrogation.trains}{EXPERIMENTAL: on real recordings it found nothing beyond its surrogate null (docs/prompts/blocks/reports/W1-evaluation.md). Groups events into trains. If events come in trains, the gaps between them are of two kinds - short inside a train, long between - and a log-scale histogram of the gaps has two humps; the default rule (logISI) cuts at the valley between them, and when there is only one hump it says 'no trains' instead of inventing some. It uses Silverman's robust bandwidth; when timing inside a train is irregular it needs the gaps between trains to be \textasciitilde{}30-100x those inside, so 'no trains' can mean 'not resolvable'. Adds train\_id (-1 = none), train\_size, position\_in\_train and gap\_before\_s; the spans and their columns pass through. Tune max\_core\_gap\_s first (0 = no ceiling) if the valley lands in the wrong place; rule=adamatzky or maxinterval give comparable cuts. Put it after Event shape, and before Intervals with write\_group on to time events within trains.}
\ubblockflags{interrogation.trains}{%
  \ubflag{estimate}{}%
}
\ubblocksideinputs{interrogation.trains}{%
}
\ubblockparams{interrogation.trains}{%
  \ubparam{rule}{\texttt{logisi}}{one of \texttt{logisi}, \texttt{maxinterval}, \texttt{adamatzky}}{How a train is decided: logisi (the valley of the log-gap histogram is the cut; Pasquale et al. 2010), maxinterval (five fixed thresholds; Cotterill et al. 2016), adamatzky (a gap <= theta x the mean gap joins a word; Adamatzky 2022)}%
  \ubparam{min\_\allowbreak{}events}{3}{whole number from 2 to 1000}{Fewest events a train may hold (Pasquale and Cotterill: 3)}%
  \ubparam{max\_\allowbreak{}core\_\allowbreak{}gap\_\allowbreak{}s}{0.0}{number from 0 to 1000000}{logisi: the largest gap that can be typical inside a train, in seconds (the neuronal 100 ms); 0 = auto = no ceiling (the peaks are tried from the shortest gap up). A value > 0 also makes the cut grow cores found with it}%
  \ubparam{theta\_\allowbreak{}mult}{1.0}{number from 0.01 to 100}{adamatzky: theta = this x the mean gap (the paper used 1 and 2)}%
  \ubparam{mi\_\allowbreak{}begin\_\allowbreak{}gap\_\allowbreak{}s}{60.0}{number from 0 to 1000000}{maxinterval: a train begins at a gap <= this (s). Placeholder for fungal data}%
  \ubparam{mi\_\allowbreak{}end\_\allowbreak{}gap\_\allowbreak{}s}{120.0}{number from 0 to 1000000}{maxinterval: a train continues while gaps are <= this (s). Placeholder}%
  \ubparam{mi\_\allowbreak{}min\_\allowbreak{}ibi\_\allowbreak{}s}{180.0}{number from 0 to 1000000}{maxinterval: trains closer than this (s) merge. Placeholder}%
  \ubparam{mi\_\allowbreak{}min\_\allowbreak{}duration\_\allowbreak{}s}{0.0}{number from 0 to 1000000}{maxinterval: trains shorter than this (s) are dropped. Placeholder}%
  \ubparam{write\_\allowbreak{}group}{no}{yes or no}{Write group = train id, so Intervals after this block measures timing within trains}%
  \ubparam{anchor}{\texttt{onset}}{one of \texttt{onset}, \texttt{start}}{Time each event from its measured onset (Event shape upstream) or its span start}%
}
